\chapter{Introduction}
\label{chap:introduction}

Concevoir et maintenir des systèmes logiciels de grande taille suppose de pouvoir naviguer dans une multiplicité de représentations --- code, diagrammes, documents --- sans perdre la cohérence entre elles. Les approches traditionnelles, centrées sur l'écriture manuelle du code, peinent à tenir ce contrat à mesure que la complexité s'accroît~: les artefacts dérivent, les contraintes métier se diluent dans l'implémentation et la documentation prend du retard.

C'est dans ce contexte que s'inscrit l'Ingénierie Dirigée par les Modèles (\textit{Model-Driven Engineering}, MDE). Cette méthodologie place le modèle — une représentation abstraite et simplifiée du système — au cœur du processus de développement. En se concentrant sur les concepts métier essentiels et en faisant abstraction des détails d'implémentation, le MDE permet, par le biais de transformations automatisées, de générer des artefacts tels que du code source ou de la documentation technique.

Pour soutenir cette approche, des outils de génie logiciel assisté par ordinateur, appelés outils CASE (\textit{Computer-Aided Software Engineering}), ont été développés. Ces outils souffrent cependant d'une rigidité méthodologique : ils sont limités à un ensemble fixe de notations de modélisation, codées en dur par le fournisseur. Pour pallier ce manque de flexibilité, les outils MetaCASE, tels que MetaDONE, ont émergé. Ils permettent de définir des langages de modélisation spécifiques à un domaine (\textit{Domain Specific Modeling Languages}, DSML) afin de générer dynamiquement des environnements de modélisation sur mesure.

Si ces outils offrent des gains substantiels en termes de réutilisation et de standardisation, ils font face à un obstacle technique majeur : la scalabilité. Lorsque la taille des modèles augmente, les performances des outils MDE peuvent se dégrader significativement, en particulier lors des phases de persistance et de chargement en mémoire~\cite{kolovos2013research}.

Ce mémoire se concentre sur l'analyse empirique et l'optimisation des performances de MetaDONE, un outil MetaCASE développé à l'Université de Namur~\cite{Eng2007}. MetaDONE repose sur le langage de métamodélisation MetaL, organisé en deux couches complémentaires : MetaL$_1$, le noyau formel qui définit les concepts atomiques de persistance, et MetaL$_2$, la couche d'abstraction qui permet de construire et de manipuler des métamodèles spécifiques à un domaine. À ces deux couches s'ajoute GraSyLa, un DSML déclaratif permettant de spécifier les notations concrètes (graphiques ou textuelles) associées à un métamodèle, dont l'évaluation conditionne directement la réactivité de l'interface utilisateur. La propriété fondamentale qui distingue MetaL de ses concurrents — et qui en fait à la fois la force et la source de défis techniques — est sa réflexivité totale : le langage se décrit lui-même avec ses propres construits, formant un graphe de données dense et cyclique plutôt qu'une arborescence. Cette propriété est indispensable au \textit{bootstrapping} de l'outil, qui peut ainsi s'étendre sans modification du code source, mais elle engendre des surcoûts spécifiques lors de la persistance et du chargement des modèles.

MetaDONE est développé en Java et s'appuie sur l'architecture modulaire OSGi (\textit{Open Services Gateway initiative}), qui permet d'étendre l'outil par des plugins indépendants — support de nouvelles notations (BPMN, réseaux de Petri, diagrammes entité-association), nouveaux mécanismes de persistance — sans modifier le noyau de l'application. Cette extensibilité constitue une contrainte autant qu'un avantage : toute optimisation du noyau doit rester compatible avec l'ensemble des plugins existants, ce qui renforce la nécessité d'une méthodologie de profilage rigoureuse et d'optimisations locales et non invasives.

La question centrale de ce mémoire est la suivante : dans quelle mesure les surcoûts liés à la réflexivité et à la persistance de MetaDONE sont-ils mesurables et réductibles, sans remettre en cause les garanties architecturales de l'outil ?

Pour y répondre, nous adoptons une démarche empirique structurée autour de scénarios représentatifs des usages réels de MetaDONE, en appliquant un cycle itératif mesure--analyse--optimisation. Les outils de profilage Java retenus --- Java Flight Recorder pour les mesures globales à faible surcoût, et async-profiler pour l'analyse fine des méthodes chaudes par flamegraphs --- permettent de localiser avec précision les goulots d'étranglement sans biais d'observation significatif.

Les principales contributions de ce travail sont les suivantes. Premièrement, l'identification empirique de quatre goulots d'étranglement caractéristiques de MetaDONE, situés respectivement dans la construction de métamodèles, le chargement de scripts GraSyLa, la boucle de visualisation et la création de modèles concrets à grande échelle. Deuxièmement, l'implémentation et la validation expérimentale de deux optimisations algorithmiques compatibles avec le Design by Contract de l'outil~: la mise en cache de la vérification d'appartenance dans \texttt{ConcreteModel\_impl} (ramenant la construction de métamodèles de $O(N^2)$ à $O(N)$) et la correction d'une asymétrie de conception dans \texttt{\_createRole} (ramenant la création de rôles de $O(N^2)$ à un plancher proche de $O(N \log N)$, soit un gain médian de $4{,}39\times$ mesuré à $N = 5\,000$ sur 10 répétitions par campagne). Troisièmement, le rejet motivé de deux candidates d'optimisation~: l'une dont les mesures montrent qu'elle n'a aucun effet sur le backend utilisé, l'autre dont l'implémentation violerait le principe Design by Contract de l'outil. Enfin, un profilage live de l'application Eclipse qui quantifie le coût réel du pipeline de visualisation et oriente les perspectives d'amélioration.

Ce mémoire s'organise comme suit. Le chapitre~\ref{chap:etat_art} établit l'état de l'art du MDE, des outils MetaCASE et des problématiques de scalabilité et de persistance de modèles. Le chapitre~\ref{chap:architecture} présente l'architecture interne de MetaDONE dans laquelle s'inscrivent les optimisations. Le chapitre~\ref{chap:methodologie} décrit la méthodologie expérimentale adoptée. Le chapitre~\ref{chap:developpement} constitue le cœur du travail : il présente, pour chaque scénario, le diagnostic des goulots d'étranglement, les optimisations implémentées et leur évaluation quantitative. Le chapitre~\ref{chap:discussion} discute les résultats en les replaçant dans le contexte architectural de MetaDONE et identifie les perspectives d'amélioration. Le chapitre~\ref{chap:conclusion} conclut le mémoire.